% BEGIN THERMAL_R03_SYNC_18
\section{Évaluation des constantes et comparaison quantitative}
\label{sec:ii-contraste-cuantitativo}

Les constructions précédentes permettent de séparer trois opérations :
produire une grandeur, exprimer sa coordonnée dans une coordinatisation et comparer cette
coordonnée à une référence. L’état APP--TRIT--TPK conserve les feuillets,
l’orientation, le quotient, la retenue et la mémoire avant ces lectures. Son
prolongement nonadique détermine les coordonnées corrélées
\(\pi,\varphi,e,\alpha\) ; la norme multiplicative sur le quotient local et le retour régional
produisent les sections d’action ; l’incidence marquée \(90/120\) et
les canaux conjugués produisent la fonctionnelle de Barbero ; le transport
angulaire détermine ensuite les coordonnées CKM. Aucune entrée des tableaux
externes n’intervient dans le choix de ces opérations.

On réunit ici les évaluations des constructions précédentes et leurs
références de comparaison. Les chiffres des sorties internes expriment
une précision de calcul, non une incertitude expérimentale. Pour une valeur
\(x\) et une référence non nulle \(x_0\), toutes deux dans la même coordinatisation, on utilise
\[
 \Delta x=x-x_0,\qquad r_x=\frac{x-x_0}{x_0}.
\]
Lorsque la référence possède des incertitudes marginales asymétriques
\(u_-,u_+\), nous définissons en outre le résidu descriptif
\begin{equation}
 z_x=\begin{cases}
 \Delta x/u_+,&\Delta x\geq0,\\
 \Delta x/u_-,&\Delta x<0.
 \end{cases}
 \label{eq:ii-residuo-marginal}
\end{equation}
Cette normalisation ne présuppose pas l’indépendance entre les lignes et ne constitue pas
à elle seule un test statistique conjoint.

\subsection{Constante réduite de Planck et action par tour : deux orientations}

Dans la coordinatisation \(\mathcal U_S\mapsto\mathrm{J\,s}\), les sections déjà
construites sont
\begin{align}
 \hbar_{\rm pre}&=H_5\,10^{-34}\,\mathrm{J\,s},&
 H_5&=1.0545718176461544696\ldots,\\
 \hbar_{\rm ret}&=\eta_{\rm ret}\,10^{-34}\,\mathrm{J\,s},&
 \eta_{\rm ret}&=1.0545718169150390611\ldots,\\
 h_a&=2\pi\hbar_a,&a&\in\{\mathrm{pre},\mathrm{ret}\}.
 \label{eq:ii-cuatro-coordenadas-planck}
\end{align}
La décade provient de \(\Sigma_H=\varphi\alpha^{16}\) ; la différence entre
les deux sections provient du retour régional. Multiplier par \(2\pi\)
fait passer de l’action par radian à l’action par tour, sans introduire une autre
orientation ni altérer le résidu relatif.

La référence du SI est
\[
 h_{\rm SI}=6.62607015\times10^{-34}\,\mathrm{J\,s},\qquad
 \hbar_{\rm SI}=\frac{h_{\rm SI}}{2\pi}
 =1.0545718176461563912\ldots\times10^{-34}\,\mathrm{J\,s}.
\]
Les deux expressions sont exactes. Le développement infini de \(\hbar_{\rm SI}\)
n’est pas une incertitude de mesure. La fixation de la valeur numérique de
\(h\) appartient à la définition du SI en vigueur depuis 2019
\cite{bipm2018} ; la comparaison s’effectue après la construction HMT.

\begin{table}[htbp]
\centering\small
\setlength{\tabcolsep}{5pt}
\begin{tabular}{@{}llr@{}}
\toprule
Grandeur & Coordonnée en \(10^{-34}\,\mathrm{J\,s}\) & \(r_x\) par rapport au SI\\
\midrule
\(\hbar_{\rm pre}\) & \(1.0545718176461544696\) & \(-1.822221\times10^{-15}\)\\
\(\hbar_{\rm ret}\) & \(1.0545718169150390611\) & \(-6.932836\times10^{-10}\)\\
\(h_{\rm pre}\) & \(6.62607014999998793\) & \(-1.822221\times10^{-15}\)\\
\(h_{\rm ret}\) & \(6.62607014540625433\) & \(-6.932836\times10^{-10}\)\\
\bottomrule
\end{tabular}
\caption{Action avant et après le retour. Les derniers chiffres sont
des évaluations des coordonnées obtenues ; les résidus ne sont pas des erreurs
expérimentales de \(h\) ni de \(\hbar\).}
\label{tab:ii-planck-comparacion}
\end{table}

\begin{proposition}[Invariance de la comparaison sous le tour complet]
Dans une coordinatisation commune,
\[
 \frac{h_a-h_{\rm SI}}{h_{\rm SI}}
 =\frac{\hbar_a-\hbar_{\rm SI}}{\hbar_{\rm SI}}.
\]
\end{proposition}
\begin{proof}
Les deux différences du premier membre contiennent le facteur commun
\(2\pi\), qui s’annule avec celui du dénominateur.
\end{proof}
La séparation pre/ret est quantitativement pertinente : remplacer l’une par
l’autre change la comparaison de plusieurs ordres de grandeur. Le correcteur
régional doit rester visible et ne pas être confondu avec un arrondi numérique.

\subsection{Barbero--Immirzi : valeur interne et comparaison entre constructions}

L’évaluation de la fonctionnelle se reconstruit avec ses trois contributions :
\begin{align*}
 \Delta S_{90}&=2.95169439297457621139847987861\times10^{-4},\\
 \Delta S_{120}&=1.96835112502951720093738706217\times10^{-5},\\
 \pi AC^*/180&=0.270485322934140078465586458790\ldots,\\
 \gamma_{\rm HMT}&=\pi AC^*/180+12\Delta S_{90}-\Delta S_{120}\\
 &=0.27400767269445927474725526077398041940\ldots.
\end{align*}
Cette composition conserve le poids orienté \(12:-1\), antérieur à son
évaluation décimale. Les termes indiqués reconstruisent la coordonnée
finale avec un résidu absolu inférieur à \(3\times10^{-31}\).

Comme comparaison bibliographique, on considère le dénombrement de Ghosh et Mitra,
non une mesure du paramètre. Leur équation (7), écrite avec
\(\lambda=2\pi\gamma_{\rm GM}\), est
\begin{equation}
 \sum_{j\in\{1/2,1,3/2,\ldots\}}(2j+1)
 e^{-2\pi\gamma_{\rm GM}\sqrt{j(j+1)}}=1.
 \label{eq:ii-gm-comparador}
\end{equation}
La référence est A.~Ghosh et P.~Mitra,
\emph{An improved estimate of black hole entropy in the quantum geometry
approach}, \emph{Physics Letters B} \textbf{616} (2005), 114--117,
\href{https://arxiv.org/abs/gr-qc/0411035}{arXiv:gr-qc/0411035},
équations (7)--(8). La racine est évaluée ici pour cette équation précise ;
on n’identifie pas le dénombrement des ponctures au générateur HMT.

\begin{table}[htbp]
\centering\small
\begin{tabular}{@{}lll@{}}
\toprule
Construction & Valeur adimensionnelle & Nature\\
\midrule
Fonctionnelle \(\Gamma_{6\to8}\) & \(0.2740076726944593\) & Incidence et retour HMT\\
Racine de \eqref{eq:ii-gm-comparador} & \(0.2740668582328300\) & Dénombrement bibliographique déclaré\\
\midrule
\(\gamma_{\rm HMT}-\gamma_{\rm GM}\) & \(-5.918553837\times10^{-5}\) & Différence entre constructions\\
\(\gamma_{\rm HMT}/\gamma_{\rm GM}-1\) & \(-2.159529202\times10^{-4}\) & Diferencia relativa\\
\bottomrule
\end{tabular}
\caption{Comparaison de Barbero--Immirzi dans l’orientation positive. Aucune
incertitude expérimentale ni signification statistique n’est assignée à ce tableau.}
\label{tab:ii-barbero-comparacion}
\end{table}

L’évaluation de la racine externe est reproductible sans tronquer
arbitrairement sa queue. Avec \(n=2j\), \(N=128\) et
\(r=e^{-0.27\pi}\), pour \(\gamma\geq0.27\),
\[
 \sum_{n>N}(n+1)e^{-\pi\gamma\sqrt{n(n+2)}}
 \leq \frac{r^{N+1}[(N+2)-(N+1)r]}{(1-r)^2}<10^{-43}.
\]
Chaque terme décroît strictement avec \(\gamma\). Les extrémités
\(0.27\) et \(0.28\) encadrent la racine ; la dichotomie et la borne de queue
contrôlent les décimales du tableau. Il s’agit d’une vérification du
comparateur bibliographique, non d’une voie pour sélectionner \(\gamma_{\rm HMT}\).

\subsection{Angles CKM, amplitudes et violation CP}

Le lecteur angulaire déjà construit détermine, en degrés,
\begin{equation}
 (\theta_{12},\theta_{23},\theta_{13},\delta)
 =(13.003313589509,\;2.398056157332,\;
 0.213977382244,\;65.676173123554)^\circ.
 \label{eq:ii-ckm-decimales}
\end{equation}
Ces valeurs proviennent de \(M_{\rm CKM}(A,C^*/6,\Delta)^T\) et
\(\delta=9A\). La conversion en radians s’effectue une seule fois avant
d’évaluer les sinus et les cosinus. La paramétrisation unitaire de la section
angulaire donne
\[
 |V_{\rm CKM}|\simeq
 \begin{pmatrix}
 0.974350259&0.225005836&0.003734601\\
 0.224873110&0.973489279&0.041841465\\
 0.008582400&0.041121745&0.999117283
 \end{pmatrix},
\]
et le quartet invariant conserve la phase :
\[
 J=c_{12}c_{23}c_{13}^{\,2}s_{12}s_{23}s_{13}\sin\delta
 =3.1189723326632974\times10^{-5}\quad
 \text{dans l’évaluation du lecteur angulaire.}
\]
Cette évaluation utilise \eqref{eq:ii-jarlskog} avec les coordonnées de
\eqref{eq:ii-ckm-decimales}. Les neuf modules imprimés sont arrondis ; l’unitarité appartient
à la matrice complexe avant l’arrondi, non à une matrice formée
uniquement de modules.

Le comparateur est l’ajustement à trois générations du
\href{https://pdg.lbl.gov/2025/reviews/rpp2025-rev-ckm-matrix.pdf}
{Particle Data Group, mise à jour de 2025, \emph{CKM Quark-Mixing Matrix}},
équations (12.27)--(12.28) et la valeur de \(J\) adjacente. Pour éviter
d’altérer leurs incertitudes par une transformation non linéaire, le tableau
conserve les observables exprimées par le PDG : \(s_{ij}=\sin\theta_{ij}\),
\(\delta\) en radians et \(J\).

\begin{table}[htbp]
\centering\small
\setlength{\tabcolsep}{4pt}
\begin{tabular}{@{}lrrr@{}}
\toprule
Observable & Évaluation HMT & PDG 2025 & \(z_x\)\\
\midrule
\(s_{12}\) & \(0.225007404757\) & \(0.22501\pm0.00068\) & \(-0.003817\)\\
\(s_{23}\) & \(0.041841757010\) & \(0.04183^{+0.00079}_{-0.00069}\) & \(0.014882\)\\
\(s_{13}\) & \(0.003734601164\) & \(0.003732^{+0.000090}_{-0.000085}\) & \(0.028902\)\\
\(\delta\) (rad) & \(1.146265461116\) & \(1.147\pm0.026\) & \(-0.028251\)\\
\(10^5J\) & \(3.118972333\) & \(3.12^{+0.13}_{-0.12}\) & \(-0.008564\)\\
\bottomrule
\end{tabular}
\caption{Comparaison marginale dans la paramétrisation et l’édition déclarées.
Le facteur \(10^5\) multiplie conjointement la valeur et l’incertitude de la
dernière ligne ; il ne modifie pas son résidu normalisé.}
\label{tab:ii-ckm-comparacion}
\end{table}

Les résidus sont calculés avec les coordonnées avant l’arrondi
du tableau. L’ajustement PDG combine des données et des contraintes d’unitarité ;
ses lignes ne sont pas des mesures indépendantes. Par conséquent, on ne somme pas
\(\sum z_x^2\) et on n’assigne pas de nombre de degrés de liberté à ce tableau.
Le résultat est une comparaison marginale reproductible avec cette édition,
non une mesure de probabilité du modèle. Les écarts marginaux
historiques des calculs antérieurs ne sont pas transférés à ce tableau :
leur documentation n’identifie ni l’édition ni la covariance de cette comparaison.

Le secteur PMNS maintient le transport \(U=F_\ell^\dagger F_\nu\), les
domaines spectraux et les données de phase développés dans la section
angulaire. Une ligne numérique nécessite de spécifier les repères du registre
employé et la réalisation comparée. Ce tableau n’attribue pas de chiffres CKM
au lecteur leptonique et ne détermine pas de nouvelle matrice PMNS.

\subsection{Boltzmann et gravitation universelle : grandeur, coordinatisation et référence}

La construction thermique déjà retrouvée détermine le transducteur positif
\(k_B:\mathcal L_\Theta\to\mathcal L_E\) et la relation
\begin{equation}
 k_B\Theta_{{\rm clk},a}\log3
 =\frac{h_a}{108t_0}=\frac{\hbar_a}{t_P},
 \qquad t_P=\frac{54}{\pi}t_0.
 \label{eq:ii-kb-comparacion}
\end{equation}
L’information locale, l’action et la période apparaissent avant la coordinatisation
kelvin--joule. Pour des bases positives \(u_E,u_\Theta\), si
\(k_B(u_\Theta)=b\,u_E\),
\(\Theta_{{\rm clk},a}=\theta_a u_\Theta\) et
\(E_{P,a}=\epsilon_a u_E\), l’équation équivaut à
\(b\theta_a=\epsilon_a/\log3\). Un changement
\(u'_E=s_Eu_E\), \(u'_\Theta=s_\Theta u_\Theta\) produit
\[
 b'=\frac{s_\Theta}{s_E}b,\qquad
 \theta'_a=\frac{\theta_a}{s_\Theta},\qquad
 \epsilon'_a=\frac{\epsilon_a}{s_E}.
\]
La réalisation marquée de \ref{thm:ii-ter27-transducer} détermine
en outre le coefficient au moyen de la référence spectrale et de la paire
information par référence--énergie conditionnée :
\begin{equation}
 b_*=10^{-23}\left(1+380\,10^{-3}+649\,10^{-6}\right)
     =\frac{1380649}{10^{29}}=1.380649\times10^{-23}.
 \label{eq:ii-b-evaluacion-comparada}
\end{equation}
Les degrés, leurs multiplicités et l’origine de capacité sont conservés
dans \ref{sec:ii-ter27-antecedentes} ; la caractérisation additive et
l’instrument conservatif sont développés dans
\ref{sec:ii-aditiva-instrumento}. Avec $H$ et la référence fixes,
$d\mathcal S/d\mathcal E=b_*\beta$ produit
$\Theta_P=(b_*\beta)^{-1}$ et $B_*(u_\Theta)=b_*u_E$.
Par conséquent, le tableau montre le coefficient effectivement évalué,
avec la référence SI
\[
 k_B^{\rm SI}=1.380649\times10^{-23}\,\mathrm{J\,K^{-1}},
\]
Le nombre coïncide exactement dans cette réalisation et dans les bases
transportées indiquées. Sa preuve thermométrique comprend tant la mesure
spectrale que les deux fonctionnelles : ne changer que sa normalisation change
le lecteur. L’identification de $u_E/u_\Theta$ à $\mathrm{J/K}$ est
le transport dimensionnel qui permet la comparaison. La valeur SI est
définitionnelle ; aucune incertitude expérimentale ne lui est donc assignée.

En gravité, la réalisation circulaire fixe
\begin{equation}
 G_a=\vartheta_{108}^{\circ}\frac{c_{\rm int}^{5}t_0^2}{\hbar_a},
 \qquad
 \vartheta_{108}^{\circ}=(54/\pi)^2
 =295.4525715010569\ldots.
 \label{eq:ii-g-comparacion}
\end{equation}
Le sélecteur et la réciprocité \(G_a\hbar_a=
\vartheta_{108}^{\circ}c_{\rm int}^{5}t_0^2\) sont des quantités internes à
la réalisation déclarée. Une évaluation dimensionnelle de \(G_a\) doit
conserver, outre l’orientation \(a\), les coordinatisations du temps,
de la vitesse et de l’action. La composition longitudinale et radiale de la
section~\ref{g26:sec:rigidez-radial-es} détermine d’abord
\[
 L=\frac{5759}{23040}\alpha_{\HMT}^{16}L_*,\qquad
 \hbar_{\rm ret}=\eta_{\rm ret}S_*,\qquad
 t_0=\frac{\pi_{\HMT}L}{54c_{\rm int}}.
\]
Avec la même section de vitesse constitutive, la substitution dans
\eqref{eq:ii-g-comparacion} donne
\[
 G_{\rm ret}=\frac{c_{\rm int}^3L_*^2}{S_*}
 \left(\frac{5759}{23040}\right)^2
 \frac{\alpha_{\HMT}^{32}}{\eta_{\rm ret}}.
\]
La vitesse dimensionnelle conserve la normalisation constitutive des
canaux angulaires. Avec les coordonnées déjà engendrées \(A_{\rm deg}\)
et \(C^*_{\rm deg}\), on écrit
\begin{equation}
 \begin{aligned}
 q_\pm&=\exp\!\left[-\frac{\pi_{\HMT}}{180}
               (A_{\rm deg}\pm C^*_{\rm deg})\right],\\
 r_\pm&=\frac{1-q_\pm^{90}}{1-q_\pm^{120}},\qquad
 \widehat c=\frac1{r_+r_-},\qquad c_{\rm int}=c_*\widehat c.
 \end{aligned}
 \label{eq:ii-g26-velocidad-constitutiva}
\end{equation}
Dans des bases positives compatibles,
\(\varepsilon=\varepsilon_*r_+^2\), \(\mu=\mu_*r_-^2\) et
\(c_*=(\mu_*\varepsilon_*)^{-1/2}\) ; par multiplication,
\(c_{\rm int}=(\mu\varepsilon)^{-1/2}\). Ainsi,
\begin{equation}
 G_{\rm ret}=\frac{c_*^3L_*^2}{S_*}
 r_N^2\alpha_{\HMT}^{32}\frac{\widehat c^{\,3}}{\eta_{\rm ret}}
 =\frac{L^2}{\hbar_{\rm ret}(\mu\varepsilon)^{3/2}}.
 \label{eq:ii-g26-gravedad-constitutiva}
\end{equation}
L’évaluateur utilise la section dimensionnelle \(c_{\rm int}\)
déclarée ci-après et restitue \(c_*=c_{\rm int}/\widehat c\).
La substitution vérifie l’équivalence de ces expressions dans la même
représentation : \(\widehat c\) intervient une seule fois lors du passage de
\(c_*\) à \(c_{\rm int}\) \cite{hmtg26counter}.

Dans la représentation déclarée \(L_*=1\,\mathrm m\),
\(S_*=10^{-34}\,\mathrm{J\,s}\) et
\(c_{\rm int}=299792458\,\mathrm{m\,s^{-1}}\), on obtient
\[
 L=1.616254982625126330168262501334562\ldots\times10^{-35}
 \,\mathrm m,
\]
\[
 G_{\rm ret}=6.674299659414022706367776189\ldots\times10^{-11}
 \,\mathrm{m^3\,kg^{-1}\,s^{-2}}.
\]
La valeur de comparaison \(6.67430(15)\times10^{-11}\) est conservée
séparément. La différence est
\(-3.40585977\ldots\times10^{-18}\,\mathrm{m^3\,kg^{-1}\,s^{-2}}\),
soit environ \(-0.00227057\) incertitudes-types de cette référence.
Les chiffres stables du calcul décrivent l’évaluation de la composition
avec ses lecteurs et bases explicites ; l’incertitude expérimentale appartient
à la référence mesurée.

\begin{table}[htbp]
\centering\small
\begin{tabularx}{\linewidth}{@{}lXX@{}}
\toprule
Objeto & Évaluation interne & Referencia dimensional\\
\midrule
\(B_*\) & \(1.380649\times10^{-23}\,u_E/u_\Theta\),
par \eqref{eq:ii-b-evaluacion-comparada} et la règle thermométrique marquée.
& \(1.380649\times10^{-23}\,\mathrm{J/K}\), exact dans le SI.\\[4pt]
\(G_{\rm ret}\) & \(6.6742996594140227\ldots\times10^{-11}\,
\mathrm{m^3/(kg\,s^2)}\), par \(c_{\rm int}^3L^2/\hbar_{\rm ret}\).
& \(6.67430(15)\times10^{-11}\,\mathrm{m^3/(kg\,s^2)}\), CODATA 2022.\\
\bottomrule
\end{tabularx}
\caption{Évaluations et références dans les réalisations spécifiées.
La ligne thermique conserve l’identification des bases expliquée dans le texte.
L’incertitude
relative type de la référence de \(G\) est d’environ
\(2.24743\times10^{-5}\) ; \(k_B\) est dépourvu d’incertitude définitionnelle.
L’évaluation gravitationnelle et sa différence par rapport à la référence
utilisent la composition longitudinale et radiale déclarée dans le texte.}
\label{tab:ii-kb-g-cartas}
\end{table}

Les sources externes de ce dernier tableau sont les
\href{https://www.bipm.org/en/measurement-units/si-defining-constants}
{constantes définissant le SI du BIPM} et le
\href{https://physics.nist.gov/cuu/Constants/Table/allascii.txt}
{tableau CODATA 2022 du NIST} \cite{codata2022}.
L’incertitude expérimentale de \(G\) ne se transfère ni à \(h\),
ni à \(\hbar\), ni à \(k_B\). Une égalité de chiffres dans une coordinatisation
n’élimine pas non plus la nécessité d’identifier le lecteur qui produit la grandeur.

\subsection{Reproduction des évaluations}

Le programme \texttt{verificar\_comparaciones\_II.py}, inclus avec les
sources, reçoit les coordonnées internes stockées dans des fichiers distincts
des données externes. Il vérifie la composition de la fonctionnelle de Barbero,
la conversion entre \(h\) et \(\hbar\), la racine bibliographique et sa queue,
les modules CKM, le quartet \(J\) et les résidus de ces comparaisons. Il utilise
une arithmétique décimale à 75 chiffres et des fonctions élémentaires évaluées après
le gel des sorties internes. Les instructions d’exécution
et les versions des programmes sont incluses dans le paquet reproductible.
La précision de travail n’augmente pas les chiffres disponibles du
calcul d’origine : seuls sont présentés les chiffres stables vis-à-vis de leur arrondi.
Le contrôle enregistre les fichiers utilisés et leurs empreintes, de sorte
qu’un changement d’évaluation ou d’édition externe reste visible.

L’évaluation gravitationnelle ajoutée se reproduit au moyen de
\texttt{verificar\_reproduccion.py}, dans le répertoire
\texttt{reproduccion\_gravitatoria\_20260926} des sources.
L’ensemble focal conserve les preuves de transport polaire, de composition
métrique, de rigidité simultanée, d’action de phase et de composition conjointe,
avec les substitutions structurelle et angulaire. Les vérificateurs
de substitution lisent les développements régionaux archivés et le registre
\(K\), reconstruisent les expressions ultérieures et comparent deux précisions
de travail ; les évaluations gravitationnelles conservent 71 chiffres
significatifs stables. Les identités générales sont démontrées dans le
corps du texte ; les contrôles rationnels examinent des cas finis de ces identités.
Les entrées, la provenance des copies, leurs empreintes et la portée de
chaque programme sont documentées dans le même répertoire. Cette exécution
réutilise les sorties engendrées antérieures et conserve distinctes la
reproduction focale, l’exécution intégrale du TPK et la formalisation Lean.

% NUMCHECK hbar_SI_coefficient=1.0545718176461563912 tol=1e-19
% NUMCHECK h_pre_coefficient=6.62607014999998793 tol=6e-18
% NUMCHECK h_ret_coefficient=6.62607014540625433 tol=6e-18
% NUMCHECK relative_pre=-1.822221e-15 tol=6e-22
% NUMCHECK relative_ret=-6.932836e-10 tol=6e-17
% NUMCHECK gamma_GM=0.2740668582328300 tol=6e-17
% NUMCHECK gamma_difference=-5.918553837e-5 tol=6e-15
% NUMCHECK gamma_relative=-2.159529202e-4 tol=6e-14
% NUMCHECK ckm.s12.value=0.225007404757 tol=6e-13
% NUMCHECK ckm.s23.value=0.041841757010 tol=6e-13
% NUMCHECK ckm.s13.value=0.003734601164 tol=6e-13
% NUMCHECK ckm.delta_rad.value=1.146265461116 tol=6e-13
% NUMCHECK ckm.s12.marginal_residual=-0.003817 tol=6e-7
% NUMCHECK ckm.s23.marginal_residual=0.014882 tol=6e-7
% NUMCHECK ckm.s13.marginal_residual=0.028902 tol=6e-7
% NUMCHECK ckm.delta_rad.marginal_residual=-0.028251 tol=6e-7
% NUMCHECK ckm.J.marginal_residual=-0.008564 tol=6e-7
% NUMCHECK gravity_selector=295.4525715010569 tol=6e-14
% NUMCHECK G_relative_reference_uncertainty=2.24743e-5 tol=6e-11
% END THERMAL_R03_SYNC_18
